% Standalone problem document, generated from the adjacent README.md.
% Compile with XeLaTeX. Mathematical content is maintained in Markdown.
\documentclass[11pt,a4paper]{article}
\usepackage[fontsize=10.5pt]{scrextend}
\usepackage[margin=22mm,headheight=15pt,headsep=9mm,footskip=11mm]{geometry}
\usepackage{fontspec}
\setmainfont{texgyrepagella}[Extension=.otf,UprightFont=*-regular,BoldFont=*-bold,ItalicFont=*-italic,BoldItalicFont=*-bolditalic]
\setsansfont{texgyreheros}[Extension=.otf,UprightFont=*-regular,BoldFont=*-bold,ItalicFont=*-italic,BoldItalicFont=*-bolditalic]
\setmonofont{texgyrecursor}[Extension=.otf,UprightFont=*-regular,BoldFont=*-bold,ItalicFont=*-italic,BoldItalicFont=*-bolditalic,Scale=MatchLowercase]
\usepackage{amsmath,amssymb}
\usepackage{unicode-math}
\setmathfont{texgyrepagella-math.otf}
\usepackage{xcolor,microtype,enumitem,needspace,fancyhdr}
\usepackage{bookmark}
\definecolor{ink}{HTML}{17344B}
\definecolor{accent}{HTML}{197D80}
\definecolor{muted}{HTML}{596773}
\hypersetup{colorlinks=true,linkcolor=accent,urlcolor=accent,pdftitle={IE-21 - Sharp
row-deletion singular-value limit for spherical random
matrices},pdfauthor={Open Problems in Numerical Linear Algebra}}
\urlstyle{same}
\setlength{\parindent}{0pt}
\setlength{\parskip}{5pt plus 1pt minus 1pt}
\linespread{1.04}
\setlength{\emergencystretch}{3em}
\widowpenalty=10000
\clubpenalty=10000
\displaywidowpenalty=10000
\allowdisplaybreaks[1]
\setlist{nosep,leftmargin=1.5em}
\providecommand{\tightlist}{\setlength{\itemsep}{0pt}\setlength{\parskip}{0pt}}
\providecommand{\pandocbounded}[1]{#1}
\pagestyle{fancy}
\fancyhf{}
\fancyhead[L]{\sffamily\footnotesize\color{muted}OPEN PROBLEMS IN NUMERICAL LINEAR ALGEBRA}
\fancyhead[R]{\sffamily\bfseries\footnotesize\color{accent}IE-21}
\fancyfoot[L]{\parbox[b]{0.9\textwidth}{\sffamily\fontsize{8}{10}\selectfont\color{muted}Editorial ratings; a bounded search cannot certify that a problem remains unsolved.\\Verification check: 2026-09-22}}
\fancyfoot[R]{\sffamily\footnotesize\color{muted}\thepage}
\renewcommand{\headrulewidth}{0.3pt}
\renewcommand{\headrule}{\hbox to\headwidth{\color{accent}\leaders\hrule height \headrulewidth\hfill}}
\makeatletter
\renewcommand{\section}{\Needspace{5\baselineskip}\@startsection{section}{1}{0pt}{12pt plus 2pt minus 2pt}{4pt}{\sffamily\bfseries\large\color{ink}}}
\renewcommand{\subsection}{\Needspace{4\baselineskip}\@startsection{subsection}{2}{0pt}{9pt plus 1pt minus 1pt}{3pt}{\sffamily\bfseries\normalsize\color{ink}}}
\makeatother
\setcounter{secnumdepth}{0}
\begin{document}
{\sffamily\small\color{muted}Linear systems and elimination\par}
\vspace{3pt}
{\raggedright\hyphenpenalty=10000\exhyphenpenalty=10000\sffamily\bfseries\fontsize{21}{24}\selectfont\color{ink}Sharp
row-deletion singular-value limit for spherical random matrices\par}
\vspace{7pt}
{\sffamily\small\textbf{Difficulty:} challenging\quad\textbf{Importance:} interesting
to the community\par}
{\sffamily\small\bfseries\color{accent}Status: Lean verified\par}
\vspace{4pt}
{\color{accent}\hrule height 0.8pt}
\vspace{6pt}
\subsection{Independently reviewed resolution -
2026-09-11}\label{independently-reviewed-resolution---2026-09-11}

\textbf{Affirmative resolution.} Theorem 1 and Sections 2-5 prove the
displayed Gaussian trimmed-second-moment limit in probability along
every sequence \(n\to\infty\) and \(m/n\to\infty\), with exactly
\(\lfloor\theta m\rfloor\) retained rows and the variational least
singular value. Explicit failure-probability and error bounds are
included; no faster aspect-ratio growth assumption is added.

\textbf{Author:} Matthew J. Colbrook, Department of Applied Mathematics
and Theoretical Physics, University of Cambridge.
\href{https://github.com/ajt60gaibb/OpenProblemsInNLA/blob/main/references/colbrook-recovered-2026-09-11/manuscripts/IE-21-22.pdf}{Complete
manuscript},
\href{https://github.com/ajt60gaibb/OpenProblemsInNLA/blob/main/references/colbrook-recovered-2026-09-11/verification/reviews/IE-21-22-review.md}{independent
proof review}, and
\href{https://github.com/ajt60gaibb/OpenProblemsInNLA/blob/main/references/colbrook-recovered-2026-09-11/README.md}{submission
record}. The supplied notes were reconstructed with substantial AI
assistance. This is independent agent verification, not external human
peer review or formal proof-assistant certification; no novelty or
priority claim is made.

The difficulty, importance and rating rationale below are historical
assessments of the original open target. Original statements, references
and dated audits are preserved.

\textbf{Rating rationale:} Challenging reflects a sharp limit for the
least singular value after adversarial row deletion; community impact
includes random matrix theory and robust iterative solvers.

\subsection{Lean proof and verification evidence -
2026-09-22}\label{lean-proof-and-verification-evidence---2026-09-22}

\textbf{The complete original target is Lean verified.} The
\href{https://github.com/ajt60gaibb/OpenProblemsInNLA/blob/1eb284b84ecc0d3c958d022b3e020be7fa111391/linear-systems-and-elimination/IE-21/lean/NLA/IE21/FiniteSize.lean}{formal
proof at revision 1eb284b} proves the stated probability limit for every
sequence with diverging dimension and aspect ratio. It includes the
explicit finite-size failure and error bounds, the exact floor
convention, arbitrary changing probability spaces, and the actual
normalized surface distribution. The Euclidean operator norm and
attained row-deletion minimum are proved to have their original
meanings, including rank-deficient and empty-retention cases.

\textbf{Lean formalization:} George Stepaniants, Department of Computing
and Mathematical Sciences, California Institute of Technology, with
substantial AI-agent assistance. \textbf{Matthew J. Colbrook} retains
attribution for the original mathematical proof. The original statement,
source attribution and dated history below remain preserved.

All
\href{https://github.com/ajt60gaibb/OpenProblemsInNLA/blob/main/linear-systems-and-elimination/IE-21/lean/Solution.lean}{23
required declarations} passed authentic LeanCert kernel-trust
assertions, the sandboxed Lean4 Comparator and kernel replay on non-root
Linux. Each transitive axiom closure contains only \textit{propext},
\textit{Classical.choice} and \textit{Quot.sound}. The exact analytic
proof uses no approximate numerical integration. The
\href{https://github.com/ajt60gaibb/OpenProblemsInNLA/blob/main/linear-systems-and-elimination/IE-21/lean/verification/linux/OPERATIONAL-REVIEW.md}{operational
evidence} records immutable inputs, authenticated dependency revisions
and successful isolation and rejection controls.

Two nonauthor agents independently reviewed the complete statements and
proofs:
\href{https://github.com/ajt60gaibb/OpenProblemsInNLA/blob/main/linear-systems-and-elimination/IE-21/lean/reviews/ie21-final-fidelity-evidence/review.md}{fidelity
review} and
\href{https://github.com/ajt60gaibb/OpenProblemsInNLA/blob/main/linear-systems-and-elimination/IE-21/lean/reviews/ie21-final-correctness-evidence/source-review.md}{correctness
review}. Their
\href{https://github.com/ajt60gaibb/OpenProblemsInNLA/blob/main/linear-systems-and-elimination/IE-21/lean/reviews/ie21-final-fidelity-evidence/completion-review.md}{fidelity
completion review} and
\href{https://github.com/ajt60gaibb/OpenProblemsInNLA/blob/main/linear-systems-and-elimination/IE-21/lean/reviews/ie21-final-correctness-evidence/completion-review.md}{correctness
completion review} separately check the Linux evidence. Publication
correspondence is reviewed in subsequent addenda. These are independent
AI reviews under the repository's adapted Tau Ceti protocol, not
external human peer review or official Tau Ceti endorsement. See the
\href{https://github.com/ajt60gaibb/OpenProblemsInNLA/blob/main/linear-systems-and-elimination/IE-21/lean/README.md}{project
guide} and
\href{https://github.com/ajt60gaibb/OpenProblemsInNLA/blob/main/linear-systems-and-elimination/IE-21/lean/formalization.yaml}{formalization
metadata} for exact scope, review history and reproduction instructions.

\newpage
\subsection{Problem statement}\label{problem-statement}

Fix \(0<\theta<1\). For \(A\in\mathbb R^{m\times n}\) define

\[s_\theta(A)=\min_{S\subseteq\{1,\ldots,m\},\ |S|=\lfloor\theta m\rfloor}
\min_{\|x\|_2=1}\|A_Sx\|_2,\]

where \(A_S\) contains the rows indexed by \(S\). Let \(a_\theta>0\)
satisfy \(\mathbb P(|G|\le a_\theta)=\theta\) for \(G\sim N(0,1)\), and
set

\[h_\theta=\frac1{\sqrt{2\pi}}\int_{-a_\theta}^{a_\theta}t^2e^{-t^2/2}\,dt.\]

For every integer sequence \((m_j,n_j)\) with \(n_j\to\infty\) and
\(m_j/n_j\to\infty\), let the rows of
\(A_j\in\mathbb R^{m_j\times n_j}\) be independent and uniformly
distributed on \(\mathbb S^{n_j-1}\). Is

\[\frac{s_\theta(A_j)^2}{\|A_j\|_2^2}\ \xrightarrow{\mathbb P}\ h_\theta?\]

This states Steinerberger's proposed asymptotic with \(\theta=q-\beta\).
Convergence in probability and the floor convention make the source's
limiting statement explicit. The source also asks for quantitative error
bounds; these belong to this problem rather than separate entries.

\subsection{Connection to numerical linear
algebra}\label{connection-to-numerical-linear-algebra}

The quantity measures the worst conditioning remaining after rows are
discarded. It controls convergence guarantees for quantile Kaczmarz
methods on corrupted linear systems.

\newpage
\subsection{References}\label{references}

\begin{enumerate}
\def\labelenumi{\arabic{enumi}.}
\tightlist
\item
  S. Steinerberger, \emph{Quantile-based Random Kaczmarz for corrupted
  linear systems of equations}, Information and Inference \textbf{12}(1)
  (2023), 448--465, §2.3, equation (8) and the following paragraph.
  \href{https://doi.org/10.1093/imaiai/iaab029}{Published paper}.
  \href{https://arxiv.org/abs/2107.05554}{Author preprint}, v1 (2021),
  §2.3, equation \((\diamond)\).
\item
  E. Battaglia, J.-F. Cai, J. Chen, A. Ma, D. Needell, and T. Wu,
  \emph{Quantile Randomized Kaczmarz for Streaming Linear Systems with
  Massart Noise}, arXiv:2608.27968v1 (2026), §1.1, discussion of
  Steinerberger's random-matrix heuristic and the distinction between
  static and streaming data; §5.
  \href{https://arxiv.org/abs/2608.27968}{Preprint}.
\end{enumerate}

\subsection{Earlier status check ---
2026-09-08}\label{earlier-status-check-2026-09-08}

Checked the published formulation and the arXiv version history
(2107.05554 has only v1), then searched the problem's title, the
author's name with ``subsingular'', and combinations of
``trimmed/smallest singular value'', ``quantile'', ``Kaczmarz'',
``heuristic'', ``proof'', ``2025'', and ``2026''. The August 2026
streaming paper still describes the original estimate as a heuristic;
its fresh-sample convergence bounds do not establish this static-matrix
limit. No later proof or counterexample was located in this bounded
search. That is a literature-status assessment, not a proof of openness.

\subsection{Audit update --- 2026-09-10}\label{audit-update-2026-09-10}

Rechecked Steinerberger's
\href{https://arxiv.org/pdf/2107.05554}{primary formulation}, §2.3, and
the \href{https://arxiv.org/html/2608.27968v1}{2026 streaming
follow-up}. Searches found no proof of the static proportional-deletion
limit. Cai, Chen, Ma, and Wu's new
\href{https://doi.org/10.1137/25M1785678}{subsample-size theorem}, SIAM
J. Matrix Anal. Appl. 47 (2026), 802--823, concerns quantile estimation
during QRK, a different asymptotic quantity.
\end{document}
